\documentclass[aspectratio=169]{beamer}
\usetheme{metropolis}
\metroset{sectionpage=none, numbering=counter, progressbar=frametitle}
\usepackage{amsmath, amssymb}
\usepackage[T1]{fontenc}

\title{Domain-Partition Algorithm for Turbomachinery Blades}
\subtitle{Slide 8: Streamline Simplification --- Merging \& Splitting}
\date{\today}
\author{}
\institute{}

\begin{document}

\begin{frame}{Streamline Simplification --- Merging \& Splitting}

\vspace{0.3em}
\textbf{Angle criterion} (Xiao Eq.~9)
\[
120^\circ \le \alpha \le 240^\circ \;\text{(3-valent)}, \qquad
144^\circ \le \alpha \le 216^\circ \;\text{(5-valent)}
\]

\vspace{0.3em}
\textbf{Blend formula} (Xiao Eq.~10 / Kowalski Eq.~28)
\[
\gamma_{12}(s) = (1-s)\,\gamma_1(s) + s\,\gamma_2(s)
\]

\vspace{0.4em}
\begin{itemize}
    \item Merge nearby streamlines at singularities (angle criterion)
    \item Split at transversal crossings (T-junctions)
    \item Result: clean quad-only block boundaries
\end{itemize}

\vspace{0.3em}
\small
\textit{Code:} \texttt{StreamlineMerging}, \texttt{StreamlineIntersectionSplitter}
\quad
\textit{Lit:} Xiao et al.~2020, ``Robust topology decomposition of 2D domains'' (Alg.~2)

\note{
Speaker notes:
Streamline simplification is the post-processing step that turns a raw streamline set into clean block boundaries. First, the merging step uses the angle criterion from Xiao 2020 Algorithm 2: at a 3-valent singularity two streamlines are merged if the angle between them lies between 120 and 240 degrees; at a 5-valent singularity the allowed range tightens to 144 through 216 degrees. When merging is approved, the two streamlines are blended with linear interpolation given by gamma_12(s) = (1-s) gamma_1(s) + s gamma_2(s). Second, the splitting step detects transversal crossings and inserts T-junctions so that every block boundary is a simple polyline. The combined result is a set of quad-only block boundaries with no dangling or overlapping segments. In the code, StreamlineMerging performs the angle test and blend, while StreamlineIntersectionSplitter handles the T-junction insertion. Both are called in sequence inside tmesh_partition.py around lines 1071--1074.
}

\end{frame}

\end{document}
